\documentclass[11pt]{article}
\usepackage[margin=1in]{geometry}
\usepackage{amsmath,amssymb}
\usepackage[hidelinks]{hyperref}
\setlength{\parindent}{0pt}
\setlength{\parskip}{0.55em}
\setlength{\emergencystretch}{1em}
\title{STAT 5200: Project Proposal}
\author{Anish Ganesh}
\date{October 29, 2026}
% Generated by code/descriptives.R; do not edit numeric values manually.
\newcommand{\HWTotalN}{453}
\newcommand{\HWMainN}{359}
\newcommand{\HWControlN}{187}
\newcommand{\HWTreatedN}{172}
\newcommand{\HWPostCompleteN}{352}
\newcommand{\HWAllCompleteN}{341}
\newcommand{\HWPreTimeN}{347}
\newcommand{\HWPostTimeN}{356}
\newcommand{\HWPreGradeN}{352}
\newcommand{\HWPostGradeN}{353}
\newcommand{\HWMeanPreTime}{29.37}
\newcommand{\HWMeanPostTime}{22.24}
\newcommand{\HWMeanPreGrade}{3.88}
\newcommand{\HWMeanPostGrade}{4.12}


\begin{document}
\maketitle

\section*{Replication and Project Proposal}
\textbf{Can AI help professionals meet quality and time targets?}

\textbf{Problem.}
My work on AI assistants motivates a practical question: does access to
generative AI increase the probability of producing work that meets both
a quality benchmark and a time limit? This combines quality and speed
into an outcome that is useful for evaluating a workflow.

\textbf{Data and feasibility.}
I will use Noy and Zhang's public
\href{https://osf.io/xd7qw/}{experimental data}, already downloaded locally.
Their study randomly assigned ChatGPT access between two professional
writing tasks. The cleaned sample contains \HWTotalN{} participants; I will follow
their main analysis by restricting to the \HWMainN{} participants with linear or
convex incentives. The variables are assigned treatment, task time in
minutes, mean overall grade on a 1--7 scale, occupation, task order, and
incentive group. R descriptive statistics give a mean second-task time of
\HWMeanPostTime{} minutes (\(n=\HWPostTimeN\)) and mean grade of \HWMeanPostGrade{} (\(n=\HWPostGradeN\)).
Both second-task outcomes are observed for \HWPostCompleteN{} participants.

\textbf{Replication.}
I will reproduce the completion-time result in Figure 1A of
Shakked Noy and Whitney Zhang (2023), \emph{Experimental evidence on the
productivity effects of generative artificial intelligence},
\textit{Science} 381(6654):187--192,
\href{https://doi.org/10.1126/science.adh2586}{doi:10.1126/science.adh2586}.
This will use changes between tasks and the authors' occupation-by-task-order
and incentive controls, checking the sample and estimate against their code.

\textbf{Original analysis.}
I propose that a task meets the target when its mean evaluator rating is
at least 5 out of 7 and it is completed in 20 minutes or less.
These are my proposed benchmarks, not standards set by the original study.
I will compare the proportion meeting both targets between the randomized
groups using a linear probability model with robust standard errors.
I will examine sensitivity to the chosen targets and missing outcomes.
The analysis concerns professional writing tasks in this experiment.

\textbf{Reproducibility.}
The repository contains the replication and project plans, source
links, data preparation, R descriptive statistics, and instructions:
\href{https://github.com/anishkganesh/stat5200-ai-quality-speed}{github.com/anishkganesh/stat5200-ai-quality-speed}.
\end{document}

